% kap06.tex — Hamilton-Jacobi-teori
% Hamilton-Jacobi-kompendiet

\chapter{Hamilton-Jacobi-teori}
\label{kap:hj}

Dette kapitel er kompendiets omdrejningspunkt. Alt det foregående —
Lagrange- og Hamilton-formalismen, Poisson-parenteser, kanoniske
transformationer — har været forberedelse til netop dette skridt:
at finde den kanoniske transformation, der gør dynamikken så simpel,
som den overhovedet kan blive.

\section{Strategien: transformer bevægelsen væk}

I bemærkningen sidst i afsnit~\ref{sec:hamiltons-princip}-diskussionen
i kapitel~\ref{kap:kanoniske-transf} så vi, at tidsudviklingen selv kan
opfattes som en kanonisk transformation, genereret af $H$. Dette
antyder en dristig idé: hvad hvis vi søger en kanonisk transformation
$(q,p)\to(Q,P)$, hvor den \emph{nye} Hamilton-funktion $K$ er identisk
nul?

Hvis $K\equiv 0$, siger Hamiltons ligninger i de nye variable
\eqref{eq:hamilton-nye-var} simpelthen
\begin{equation}
  \dot Q_i = \pdv{K}{P_i} = 0, \qquad \dot P_i = -\pdv{K}{Q_i} = 0.
\end{equation}
Både $Q_i$ og $P_i$ er altså \emph{konstante} — vi skriver
$Q_i=\beta_i$, $P_i=\alpha_i$, hvor $\alpha_i,\beta_i$ er $2n$
konstanter, bestemt af begyndelsesbetingelserne. Al den dynamiske
information om systemet er dermed flyttet over i selve
transformationen: at finde bevægelsen $q_i(t)$ er reduceret til at
finde \emph{transformationen} — dvs.\ den genererende funktion — der
opnår $K\equiv 0$. Dette er den \emph{eneste} idé i Hamilton-Jacobi-
teorien; resten af kapitlet handler om at gøre den eksplicit.

\section{Udledning af Hamilton-Jacobi-ligningen}

Vi søger transformationen via en genererende funktion af type $F_2$
(afsnit~\ref{kap:kanoniske-transf}'s notation), som vi her — af
historiske og notationsmæssige grunde — kalder $S(q,t)$ i stedet for
$F_2(q,P,t)$, og hvor vi lader $S$ afhænge af $n$ konstanter
$\alpha_1,\dots,\alpha_n$, som spiller rollen af de nye impulser $P_i$:
\begin{equation}
  S = S(q_1,\dots,q_n,\alpha_1,\dots,\alpha_n,t).
\end{equation}
Ifølge de $F_2$-relationer, vi udledte i \eqref{eq:F2-relationer}
(med $P_i\to\alpha_i$), er
\begin{equation}
  p_i = \pdv{S}{q_i}, \qquad Q_i = \pdv{S}{\alpha_i}, \qquad K = H + \pdv{S}{t}.
  \label{eq:S-relationer}
\end{equation}
Betingelsen $K\equiv 0$ giver derfor, ved indsættelse af den første
relation i $H$:
\begin{equation}
  \boxed{H\left(q_1,\dots,q_n,\pdv{S}{q_1},\dots,\pdv{S}{q_n},t\right) + \pdv{S}{t} = 0.}
  \label{eq:hj-ligning}
\end{equation}
Dette er \emph{Hamilton-Jacobi-ligningen}. Det er en enkelt
\emph{partiel} differentialligning, første ordens i tid og i hver
$q_i$, for den ukendte funktion $S(q,t)$ — i modsætning til Hamiltons
$2n$ almindelige, om end koblede, differentialligninger for
$(q_i(t),p_i(t))$. Funktionen $S$ kaldes \emph{Hamiltons
principfunktion}.

\begin{bemaerkning}
Det virker måske som en dårlig handel at bytte $2n$ almindelige
differentialligninger ud med \'en partiel differentialligning — en
PDE er typisk sværere at løse end et system af ODE'er. Værdien af
Hamilton-Jacobi-ligningen ligger ikke i, at den generelt er lettere at
løse, men i, at den for en stor og vigtig klasse af systemer
(\emph{separable} systemer, se afsnit~\ref{sec:separation}) reducerer
til simple, en-variabel integraler — og i, at den giver et begrebsapparat
(virkningen $S$ som funktion af tilstand og tid), der overlever direkte
ind i kvantemekanikken, som vi antyder i afsnit~\ref{sec:hj-og-kvante}.
\end{bemaerkning}

\subsection{$S$ er virkningen}

Vi bekræfter, at $S$ virkelig er den samme størrelse, vi kaldte
virkningen i \eqref{eq:virkning}. Langs en faktisk løsning, hvor
$p_i=\partial S/\partial q_i$ (fra \eqref{eq:S-relationer}):
\begin{equation}
  \dv{S}{t} = \sum_i\pdv{S}{q_i}\dot q_i + \pdv{S}{t}
  = \sum_i p_i\dot q_i + \pdv{S}{t}
  = \sum_i p_i\dot q_i - H = L,
\end{equation}
hvor vi i næstsidste skridt brugte Hamilton-Jacobi-ligningen
\eqref{eq:hj-ligning} til at sætte $\partial S/\partial t=-H$, og i
sidste skridt brugte definitionen af $L$ fra Legendre-transformationen
\eqref{eq:def-H}. Vi finder altså $dS/dt = L$, hvilket ved integration
giver
\begin{equation}
  S = \int L\,dt + \text{konstant},
\end{equation}
netop virkningen \eqref{eq:virkning}, evalueret langs den faktiske
bane, som funktion af dens sluttidspunkt og -position. Dette er ikke
en tilfældig sammenfald af navne: Hamilton-Jacobi-ligningen er den
naturlige PDE-formulering af, at virkningen gør sig selv stationær.

\subsection{Den fulde (komplette) løsning og bevægelsen}

En løsning $S(q,\alpha,t)$ til \eqref{eq:hj-ligning}, som afhænger af
$n$ uafhængige konstanter $\alpha_1,\dots,\alpha_n$ på en sådan måde,
at Hesse-matricen $\partial^2S/\partial q_i\partial\alpha_j$ er
ikke-singulær, kaldes en \emph{komplet integral} af
Hamilton-Jacobi-ligningen. (Bemærk: da kun \emph{afledte} af $S$
optræder i \eqref{eq:hj-ligning}, er $S$ i sig selv altid bestemt op
til en additiv konstant — denne tæller ikke som en af de $n$ ægte
$\alpha_i$'er.)

Givet en sådan komplet integral, giver den anden relation i
\eqref{eq:S-relationer} os direkte de $n$ resterende konstanter:
\begin{equation}
  \beta_i = \pdv{S}{\alpha_i}(q,\alpha,t).
  \label{eq:beta-relation}
\end{equation}
Dette er $n$ ligninger, der (i princippet) kan løses for
$q_1,\dots,q_n$ som funktion af $t$ og de $2n$ konstanter
$(\alpha,\beta)$ — og \emph{det} er selve løsningen til bevægelsen! Vi
opsummerer proceduren:

\begin{enumerate}
  \item Løs Hamilton-Jacobi-ligningen \eqref{eq:hj-ligning} for en
        komplet integral $S(q,\alpha,t)$.
  \item Beregn $\beta_i=\partial S/\partial\alpha_i$; sæt disse lig
        med konstanter, bestemt af begyndelsesbetingelserne.
  \item Løs de resulterende ligninger for $q_i(t)$.
  \item Find (om ønsket) $p_i(t)=\partial S/\partial q_i$, evalueret
        langs den fundne løsning.
\end{enumerate}

Denne procedure er hele Hamilton-Jacobi-metoden. Det eneste
resterende arbejde — som kan være betydeligt — er trin 1: at finde en
komplet integral. Vi vender os nu mod det vigtigste tilfælde, hvor
dette er muligt i praksis.

\section{Tidsuafhængigt $H$: separation i tiden}
\label{sec:separation}

Antag, at $H$ ikke afhænger eksplicit af $t$ (som i alle vores
eksempler hidtil). Vi ved fra kapitel~\ref{kap:hamilton}, at $H$ da er
en bevaret størrelse langs bevægelsen; kald denne bevarede værdi $E$
(typisk den mekaniske energi, jf.\ diskussionen efter
\eqref{eq:dHdt-generel}). Det er da naturligt at gætte på en løsning
af formen
\begin{equation}
  S(q,\alpha,t) = W(q,\alpha) - Et,
  \label{eq:S-separation-tid}
\end{equation}
hvor $E=E(\alpha)$ er en af konstanterne (eller en funktion af dem —
vi vælger simpelthen $\alpha_1\equiv E$). Indsat i
\eqref{eq:hj-ligning}, og da $\partial S/\partial q_i=\partial W/\partial q_i$
(det tidsafhængige led bidrager intet til de rumlige afledte):
\begin{equation}
  \boxed{H\left(q_1,\dots,q_n,\pdv{W}{q_1},\dots,\pdv{W}{q_n}\right) = E.}
  \label{eq:hj-tidsuafhaengig}
\end{equation}
Dette er den \emph{tidsuafhængige Hamilton-Jacobi-ligning}, og $W$
kaldes \emph{Hamiltons karakteristiske funktion}. Den er en PDE i de
rumlige variable alene — tiden er helt elimineret — hvilket typisk gør
den væsentligt lettere at løse end \eqref{eq:hj-ligning}.

\section{Eksempel: den harmoniske oscillator}

\begin{eksempel}
Vi løser den énddimensionale harmoniske oscillator
(jf.\ eksemplet i kapitel~\ref{kap:hamilton}) fuldstændigt via
Hamilton-Jacobi-metoden. Hamilton-funktionen er
\begin{equation}
  H = \frac{p^2}{2m} + \frac12 kx^2.
\end{equation}
Den tidsuafhængige Hamilton-Jacobi-ligning \eqref{eq:hj-tidsuafhaengig}
(med $p\to dW/dx$) bliver
\begin{equation}
  \frac{1}{2m}\left(\dv{W}{x}\right)^2 + \frac12 kx^2 = E
  \quad\Longrightarrow\quad
  \dv{W}{x} = \sqrt{2mE - mkx^2},
\end{equation}
hvor vi vælger den positive rod (svarende til bevægelse i positiv
$x$-retning; se bemærkningen nedenfor). Integration giver
\begin{equation}
  W(x,E) = \int\sqrt{2mE-mkx^2}\,dx,
\end{equation}
og dermed, ifølge \eqref{eq:S-separation-tid},
\begin{equation}
  S(x,E,t) = \int\sqrt{2mE-mkx^2}\,dx - Et.
\end{equation}
Vi finder $\beta$ ved \eqref{eq:beta-relation}, med $\alpha_1=E$:
\begin{equation}
  \beta = \pdv{S}{E} = \int\pdv{E}\sqrt{2mE-mkx^2}\,dx - t
        = \int\frac{m}{\sqrt{2mE-mkx^2}}\,dx - t.
\end{equation}
Skriv $a\equiv\sqrt{2E/k}$ (amplituden) og $\omega\equiv\sqrt{k/m}$
(vinkelfrekvensen, jf.\ kapitel~\ref{kap:hamilton}). Da
$2mE-mkx^2=mk(a^2-x^2)$, bliver integranden
\begin{equation}
  \frac{m}{\sqrt{mk(a^2-x^2)}} = \frac{1}{\omega}\cdot\frac{1}{\sqrt{a^2-x^2}},
\end{equation}
og standardintegralet $\int dx/\sqrt{a^2-x^2}=\arcsin(x/a)$ giver
\begin{equation}
  \beta = \frac{1}{\omega}\arcsin\left(\frac{x}{a}\right) - t.
\end{equation}
Løst for $x$:
\begin{equation}
  \boxed{x(t) = a\sin\bigl(\omega(t+\beta)\bigr),}
\end{equation}
den velkendte harmoniske bevægelse — udledt her fuldstændigt fra
Hamilton-Jacobi-ligningen, uden på noget tidspunkt at have opskrevet
eller løst Newtons eller Hamiltons differentialligninger direkte! Som
et sidste tjek: $p=\partial W/\partial x=\sqrt{2mE-mkx^2}=m a\omega\cos(\omega(t+\beta))$,
som netop stemmer med $p=m\dot x$ udregnet fra $x(t)$ ovenfor.

\begin{bemaerkning}
Fortegnsvalget for kvadratroden i $dW/dx$ svarer til, at $W$ (og dermed
$S$) kun er defineret entydigt inden for én "gren" af bevægelsen (fx
mens $x$ vokser). Ved vendepunkterne $x=\pm a$, hvor $p=0$, skal man i
princippet skifte til den anden gren. Dette tekniske punkt bliver
vigtigt, når vi i kapitel~\ref{kap:virkning-vinkel} indfører
virkningsvariable, som netop er konstrueret til at give en glat,
global beskrivelse hen over sådanne vendepunkter.
\end{bemaerkning}
\end{eksempel}

\section{Det geometriske billede: $S$ som en "bølgefront"}

Der er en frugtbar geometrisk måde at tænke på Hamilton-Jacobi-
ligningen, som binder direkte tilbage til Fermats princip fra
kapitel~\ref{kap:newton-variation}. For fastholdt $t$ definerer
$S(q,t)=\text{konstant}$ en familie af "overflader" (kurver, for
$n=1$; hyperflader for større $n$) i konfigurationsrummet — analogt til
bølgefronter i optikken. Bevægelsens impuls $p_i=\partial S/\partial q_i$
er, geometrisk set, normalvektoren til disse overflader — ganske som
lysstråler i geometrisk optik altid står vinkelret på bølgefronterne.
Partiklens faktiske bane $q_i(t)$ krydser derfor altid
$S$-overfladerne vinkelret, og bevæger sig fra overflade til overflade,
efterhånden som $t$ (og dermed værdien af $S$ langs banen) vokser.

\begin{figure}[htbp]
  \centering
  \begin{tikzpicture}[scale=1.0]
    % "Bølgefronter" = niveaukurver for S ved stigende tidspunkter
    \foreach \x/\op in {0.4/0.35,1.3/0.5,2.2/0.65,3.1/0.8,4.0/1.0} {
      \draw[thick,blue!60!black,opacity=\op]
        (\x,-1.6) to[out=90,in=-100] (\x-0.3,0) to[out=80,in=-90] (\x,1.6);
    }
    \node[blue!60!black] at (0.4,-2.0) {$S=S_1$};
    \node[blue!60!black] at (4.0,-2.0) {$S=S_5$};
    % Bane, der krydser bølgefronterne vinkelret
    \draw[-{Latex},thick,red] (0.1,-0.55) to[out=25,in=200] (1.0,0.15)
      to[out=15,in=200] (1.9,0.55) to[out=15,in=195] (2.8,0.75)
      to[out=10,in=195] (3.9,0.85);
    \node[red] at (2.2,1.15) {bane $q(t)$};
  \end{tikzpicture}
  \caption{Det geometriske billede af Hamilton-Jacobi-teorien: kurverne
  af konstant $S$ (blå, "bølgefronter", ved stigende værdier fra
  venstre mod højre) og partiklens faktiske bane (rød), som altid
  krydser dem vinkelret — impulsen $p=\partial S/\partial q$ er netop
  normalvektoren til fronterne. Sammenlign med bølgefronter og
  lysstråler i geometrisk optik, og med Fermats princip
  (kapitel~\ref{kap:newton-variation}).}
  \label{fig:hj-boelgefront}
\end{figure}

\section{Udblik: fra Hamilton-Jacobi til kvantemekanik}
\label{sec:hj-og-kvante}

Det geometriske billede i forrige afsnit er ikke kun en pædagogisk
analogi — det er historisk og begrebsmæssigt selve broen til
kvantemekanikken. I bølgeoptik er geometrisk optik (stråler vinkelret
på bølgefronter) en \emph{grænse}, gyldig når bølgelængden er meget
kort i forhold til systemets øvrige længdeskalaer; den fulde teori er
en bølgeligning for et felt, hvis faser netop svarer til $S$.

Schrödingers indsigt — som vi ikke udleder her, men blot antyder, og
som behandles fra den kvantemekaniske side i
[[kvantemekanik-kompendiet]] — var, at klassisk mekanik forholder sig
til en underliggende bølgeteori (kvantemekanikken) på præcis samme
måde, som geometrisk optik forholder sig til bølgeoptik. Skriver man
bølgefunktionen som $\psi=e^{iS/\hbar}$ og indsætter i Schrödinger-
ligningen, genfinder man i grænsen $\hbar\to0$ netop
Hamilton-Jacobi-ligningen \eqref{eq:hj-ligning} for $S$ — den klassiske
bevægelse er "korthed-af-bølgelængde"-grænsen af kvantemekanikken, med
$S/\hbar$ som bølgens fase. Dette er den dybeste grund til, at
Hamilton-Jacobi-formalismen — i modsætning til Newtons eller endda
Hamiltons ligninger direkte — er den formulering af klassisk mekanik,
der mest gennemsigtigt viser, hvor kvantemekanikken "kommer fra".

\begin{bemaerkning}
Denne forbindelse — kaldet WKB-approksimationen eller den
semiklassiske grænse — er et helt studieområde i sig selv, med
anvendelser fra tunnelering til Bohr-Sommerfeld-kvantisering.
Sidstnævnte er netop emnet for det næste kapitels afsluttende del
(kapitel~\ref{kap:adiabatisk}), hvor vi ser, hvordan virknings-
vinkel-variable (kapitel~\ref{kap:virkning-vinkel}) fører direkte til
den ældre kvanteteoris kvantiseringsbetingelser — det tankegods, der
var på plads, da kvantemekanikken blev født, og som er dette
kompendiums endestation.
\end{bemaerkning}

\section{Opsummering}

\begin{itemize}
  \item Strategien: find en kanonisk transformation til
        $K\equiv 0$, hvorved de nye variable $Q_i=\beta_i$,
        $P_i=\alpha_i$ bliver konstanter.
  \item Med genererende funktion $S(q,\alpha,t)$ af type $F_2$ giver
        $K=H+\partial S/\partial t=0$ Hamilton-Jacobi-ligningen
        $H(q,\partial S/\partial q,t)+\partial S/\partial t=0$.
  \item $S$ er (op til en konstant) selve virkningen: $dS/dt=L$.
  \item En komplet integral $S(q,\alpha,t)$ giver bevægelsen direkte
        via $\beta_i=\partial S/\partial\alpha_i$, løst for $q_i(t)$.
  \item Når $H$ er tidsuafhængig: $S=W(q,\alpha)-Et$, med $W$
        (Hamiltons karakteristiske funktion) løsning til den
        tidsuafhængige Hamilton-Jacobi-ligning
        $H(q,\partial W/\partial q)=E$.
  \item Geometrisk: $S=\text{konstant}$-flader er "bølgefronter", og
        banen krydser dem vinkelret med $p=\partial S/\partial q$ som
        normalvektor — analogt til geometrisk optik, og forløber for
        den semiklassiske grænse af kvantemekanikken ($\psi\sim e^{iS/\hbar}$).
\end{itemize}

\begin{opgave}
Løs den tidsuafhængige Hamilton-Jacobi-ligning for en fri partikel i
én dimension, $H=p^2/2m$, og vis, at proceduren giver den forventede
løsning $x(t)=x_0+\frac{p_0}{m}t$.
\end{opgave}

\begin{opgave}
Betragt en partikel, der falder frit i et konstant tyngdefelt,
$H=p^2/2m+mgx$ (med $x$ målt opad). Opstil og løs den tidsuafhængige
Hamilton-Jacobi-ligning, og genfind $x(t)=x_0+v_0t-\frac12gt^2$ ved
proceduren i afsnit~\ref{sec:separation}.
\end{opgave}

\begin{opgave}
Bekræft ved direkte differentiation, at $p=\partial W/\partial x=
\sqrt{2mE-mkx^2}$ og $x(t)=a\sin(\omega(t+\beta))$ (fra eksemplet med
den harmoniske oscillator) tilsammen opfylder $p=m\dot x$, og at
$H=p^2/(2m)+\frac12kx^2=E$ er opfyldt identisk for alle $t$.
\end{opgave}
